\par\noindent\textbf{(1)} 由互易线性磁路，磁通链 λ₁=L₁i₁+Mi₂、λ₂=Mi₁+L₂i₂，电压按被动符号约定为 v=dλ/dt= [L₁,M;M,L₂] di/dt。互感矩阵对称是能量可积与Lorentz互易的结果。
\par\noindent\textbf{(2)} 对准静态缓慢建立电流过程，外源做功增量 dU=i₁dλ₁+i₂dλ₂；在线性常系数情形积分得到 U=1/2(L₁i₁²+2Mi₁i₂+L₂i₂²)。交叉项的正负与绕组相对绕向有关，但完整二次型必须对全部i非负。
\par\noindent\textbf{(3)} 对称矩阵正半定要求L₁,L₂≥0且行列式 L₁L₂−M²≥0，故 |M|≤√(L₁L₂)。定义 k=M/√(L₁L₂) 有 |k|≤1。|k|=1 时某电流组合储能为零，理想完全耦合的漏感消失；真实线圈通常 |k|<1。
\par\noindent\textbf{(4)} 若次级短路且电压v₂=0，准静态下 di₂/dt=−(M/L₂)di₁/dt。代入 v₁=[L₁−M²/L₂]di₁/dt，故主级有效电感 L\_eq=L₁(1−k²)≥0；完全耦合时趋零，而开路次级i₂=0时主级仍显示L₁。
